% Additive canonical fragment for the Hurwitz/Stieltjes volume.
% Place after 07-mixed-twist-identities.tex. No repository files were changed.
% Requires amsmath, amssymb, amsthm and the canonical theorem/proof environments.
% All new labels use the utj: namespace; no new global macros are introduced.
% Full details, literature attribution and additional formulas are in article.tex.

\section{Unequal-twist jets: exact expansions and finite functional closure}
\label{utj:sec:main}
Retain the twisted Fourier convention
\[
 \widehat K_\alpha^\theta(n)=(2\pi i(n+\theta))^{-\alpha},\qquad
 L_{n,\theta}=\log(2\pi|n+\theta|)
       +\frac{i\pi}{2}\operatorname{sgn}(n+\theta).
\]
In this section write
$J_m^{r,\theta}=(-1)^m\partial_\alpha^mK_\alpha^\theta|_{\alpha=r}$.
Its multiplier is $L_{n,\theta}^m(2\pi i(n+\theta))^{-r}$.
Every Fourier mode is retained. In particular $K_0^\theta=\delta_0$.
The finite integer-order partial-fraction theorem does not by itself
permit differentiation in those discrete order labels. A convergent
identity in arbitrary orders supplies the required extension.

\begin{theorem}[Central expansion and every mixed order jet]
\label{utj:thm:central}
Let $\theta_j\in(0,1)$ and choose $c\in(0,1)$ with
$R=\max_j|\theta_j-c|<\eta=\min(c,1-c)$. The midpoint of the extreme
twists always satisfies this condition. Put $d_j=\theta_j-c$,
$A=\sum_j\alpha_j$, and define the polynomials $h_k$ by
\begin{align}
 \sum_{k\ge0}h_k(\boldsymbol\alpha;\boldsymbol d)t^k
   &=\prod_j(1-d_jt)^{-\alpha_j},\label{utj:eq:hgen}\\
 h_0=1,\qquad k h_k
   &=\sum_{\ell=1}^k\left(\sum_j\alpha_jd_j^\ell\right)h_{k-\ell}.
 \label{utj:eq:hrec}
\end{align}
Then, locally normally in the distribution topology,
\begin{equation}
 \mathop{*}_{j=1}^d K_{\alpha_j}^{\theta_j}
  =\sum_{k\ge0}(-2\pi i)^k h_k(\boldsymbol\alpha;\boldsymbol d)K_{A+k}^c.
 \label{utj:eq:central}
\end{equation}
All order derivatives are allowed. In particular, with
$M=\sum_jm_j$, $R_*=\sum_jr_j$, and the usual multi-index conventions,
\begin{align}
 \mathop{*}_{j=1}^d J_{m_j}^{r_j,\theta_j}
 =\sum_{k\ge0}(-2\pi i)^k\sum_{\boldsymbol q\le\boldsymbol m}
    (-1)^{|\boldsymbol q|}\binom{\boldsymbol m}{\boldsymbol q}
    (\partial^{\boldsymbol q}h_k)(\boldsymbol r;\boldsymbol d)
    J_{M-|\boldsymbol q|}^{R_*+k,c}.
 \label{utj:eq:jets}
\end{align}
\end{theorem}
\begin{proof}
At each Fourier mode,
$L_{n,\theta_j}=L_{n,c}+\log(1+d_j/(n+c))$ because the ratio is
positive. Expand the binomial factors and use
$(n+c)^{-k}=(2\pi i)^k(2\pi i(n+c))^{-k}$.
For $R<R_0<\eta$, Cauchy's estimate applied to
\eqref{utj:eq:hgen} gives $|h_k|\le C R_0^k$ on compact order sets.
The resulting modewise absolute sum has a common polynomial growth
bound in $n$. Pairing against rapidly decreasing smooth Fourier
coefficients proves normal convergence and Fourier uniqueness proves
\eqref{utj:eq:central}. Cauchy's formula in the order variables and
Leibniz's rule give \eqref{utj:eq:jets}, including its sign.
\end{proof}

All coefficients are finite Pochhammer derivatives. With unsigned
Stirling numbers $\left[\begin{smallmatrix}k\\q\end{smallmatrix}\right]$,
\[
 \left.\partial_u^q\frac{(u)_k}{k!}\right|_{u=0}
  =\frac{q!}{k!}\left[\begin{smallmatrix}k\\q\end{smallmatrix}\right].
\]
For positive integer $r$, the same derivative at $r$ is
$(r)_k/k!$ times the complete Bell polynomial in
$\mathcal H_1,-\mathcal H_2,2!\mathcal H_3,\ldots$, where
$\mathcal H_j=H_{r+k-1}^{(j)}-H_{r-1}^{(j)}$.
Setting $(0)_k=0$ before differentiation would erase nonzero terms.

For arbitrary real nonintegral twists, retain a finite exceptional
Fourier head exactly and apply the same expansion only at modes with
$|n+c|>\max_j|\theta_j-c|$. This restores all branches and gives an
entire identity without requiring a common frequency interval.

\begin{theorem}[Exact finite closure criterion in the integer ladder]
\label{utj:thm:nonclosure}
Let $a_j$ be distinct real shifts and $r_j,m_j\ge0$ integers, omitting
sites with $r_j=m_j=0$. The sequence
\[
 \prod_j(n+a_j)^{-r_j}\log^{m_j}(n+a_j)
\]
is an eventual finite constant-coefficient linear combination of
$(n+c)^{-s}\log^q(n+c)$, with arbitrary finitely many centres $c$,
$s\in\mathbb Z$ and $q\ge0$, if and only if all $m_j=0$ or at most
one site is active. Thus genuine mixed logarithmic twists cannot
in general reduce to finitely many single-twist integer-order jets.
Finite Fourier corrections and derivatives of $\delta_0$ do not
remove this obstruction.
\end{theorem}
\begin{proof}
Rational partial fractions and the one-site case give sufficiency.
An eventual sequence identity in this class is an analytic identity:
a nonzero difference has a convergent expansion
$\sum_{p\ge p_0}y^{-p}P_p(\log y)$ of bounded logarithmic degree near
infinity, and its leading nonzero polynomial cannot vanish at every
large integer.

Continue the analytic identity around the finite singularities.
Monodromy difference around $-a$ adds $2\pi i$ to $\log(y+a)$
and fixes other logarithms and all integer powers. If two sites carry
logarithms, their mixed monodromy difference is nonzero on the product
but zero on every single-centre generator. If only $a$ carries a
logarithm, of degree $m$, but a denominator is present at $b\ne a$,
write the product as $R(y)\log^m(y+a)$. Applying the $(m+1)$st
monodromy difference forces higher log degrees at $a$ on the putative
right side to vanish. Applying the $m$th difference then identifies
$R(y)$ with a Laurent polynomial in $y+a$. This is impossible because
$R$ has a pole at $-b$. These arguments also hold modulo finitely many
Fourier modes, since the eventual-sequence argument is unchanged.
\end{proof}

This is a functional statement in a specified integer-order span,
not a numerical period-independence or fixed-Hurwitz-value theorem.
Neither this criterion nor a numerical relation search resolves the
remaining $S_6/S_8$ special-value problems.

\subsection{A symmetric identity and its normalization}
For $a=c-d$, $b=c+d$, $|d|<\min(c,1-c)$, set
$\eta_k=H_{2k-1}-H_{k-1}$ and $E_m^c=J_m^{0,c}$. Then
\begin{equation}
 E_1^a*E_1^b=E_2^c-
   \sum_{k\ge1}\frac{(2\pi id)^{2k}}k
              \bigl(J_1^{2k,c}+\eta_kK_{2k}^c\bigr).
 \label{utj:eq:symmetric}
\end{equation}
Indeed, at a mode, with $u=d/(n+c)$ and $L=L_{n,c}$,
\[
 (L+\log(1-u))(L+\log(1+u))
   =L^2-\sum_{k\ge1}\frac{u^{2k}}k(L+\eta_k).
\]
For the canonical coefficients
$Z_{1+u}^\theta=-\delta_0/u+\sum_{m\ge0}(-1)^mU_m^\theta u^m/m!$,
Gamma expansion gives
$E_1^\theta=-U_0^\theta-\gamma\delta_0$ and
$E_2^\theta=2U_1^\theta+2\gamma U_0^\theta+
(\gamma^2-\zeta(2))\delta_0$. Hence
\begin{align}
 U_0^a*U_0^b={}&2U_1^c+2\gamma U_0^c-\gamma(U_0^a+U_0^b)
                 -\zeta(2)\delta_0\notag\\
 &-\sum_{k\ge1}\frac{(2\pi id)^{2k}}k
              \bigl(J_1^{2k,c}+\eta_kK_{2k}^c\bigr).
 \label{utj:eq:contact}
\end{align}
The point mass is necessary, including when $a=b=c$.

\subsection{Ordinary Gamma integrals and scalar finite parts}
Write $P_\theta=J_1^{1,\theta}$, $S_0=0$, and
$S_k=\sum_{\ell=1}^k\eta_\ell/\ell$. Dividing the preceding mode
identity by $\lambda_{n,c}^2(1-u^2)$ yields
\begin{align}
 \int_0^1P_a(t)P_b(\{x-t\})\,dt
   =\sum_{k\ge0}(2\pi id)^{2k}
       \bigl(J_2^{2k+2,c}(x)-H_kJ_1^{2k+2,c}(x)
                              -S_kK_{2k+2}^c(x)\bigr).
 \label{utj:eq:gamma}
\end{align}
This is an ordinary integral: $P_\theta\in L^2$ since its coefficients
are $L_{n,\theta}/\lambda_{n,\theta}$. The coefficient double sum on
the right is absolutely convergent, dominated after summing over $k$
by $O((1+\log^2(2+|n|))/(1+|n|)^2)$, so equality is continuous and
holds at every circular argument. For $\theta=p/q$, with
$z=e^{-2\pi ip/q}$, its factors are explicitly
\begin{equation}
 P_{p/q}(x)=e^{-2\pi ipx/q}\left\{
  \sum_{\ell=0}^{q-1}z^\ell\log\Gamma\left(\frac{x+\ell}{q}\right)
                 -\frac{\gamma+\log q}{1-z}\right\},\quad 0<x<1.
 \label{utj:eq:gammafactor}
\end{equation}
To derive this representative, use
$\Phi(z,s,x)=q^{-s}\sum_\ell z^\ell\zeta(s,(x+\ell)/q)$,
$\zeta'(0,y)=\log\Gamma(y)-\tfrac12\log(2\pi)$, and the derivative
of $1/\Gamma(\alpha)$ at $\alpha=1$.

For $a,b>0$ define the sharp-cutoff constants
\[
 C_{m,n}(a,b)=\lim_{N\to\infty}\left\{
   \sum_{j=0}^{N-1}\frac{\log^m(j+a)\log^n(j+b)}{j+a}
               -\frac{\log^{m+n+1}N}{m+n+1}\right\}.
\]
Subtracting the diagonal summand leaves an absolutely convergent
series. A unilateral binomial expansion and the Pochhammer jet at
zero therefore give, with $M=m+n$ and $|b-a|<a$,
\begin{align}
 C_{m,n}(a,b)=\gamma_M(a)+(-1)^M\sum_{k\ge1}(a-b)^k
   \sum_{q=1}^{\min(n,k)}\binom nq\frac{q!}{k!}
     \left[\begin{smallmatrix}k\\q\end{smallmatrix}\right]
                           \zeta^{(M-q)}(k+1,a).
 \label{utj:eq:cutoff}
\end{align}
All zeta superscripts here are order derivatives. For other positive
shifts, split off a finite head and translate both shifts by the same
integer until the displayed condition holds. Differentiating the
convergent difference sum and applying partial fractions gives
\begin{align}
 \partial_b C_{m,n}(a,b)
   &=\frac{n}{b-a}\bigl(C_{m,n-1}(a,b)-C_{n-1,m}(b,a)\bigr),
          &&n\ge1,\ b\ne a,\label{utj:eq:recursion}\\
 C_{0,1}(a,b)&=\gamma_1(a)+\int_a^b
                  \frac{\psi(t)-\psi(a)}{t-a}\,dt,\label{utj:eq:digamma}\\
 \int_a^b C_{0,1}(a,t)\,dt
   &=(b-a)C_{0,1}(a,b)-\log\Gamma(b)+\log\Gamma(a)
                               +(b-a)\psi(a).
 \label{utj:eq:primitive}
\end{align}
The collision singularities in these divided differences are
removable. The last equality is integration by parts.

\subsection{The unilateral family and transverse ray constants}
For positive $a_j$, choose $c>0$ with $|a_j-c|<c$. The same
coefficient polynomials, now with $d_j=a_j-c$, give
\begin{equation}
 \mathcal Z(\boldsymbol\alpha;\boldsymbol a)
 :=\sum_{n\ge0}\prod_j(n+a_j)^{-\alpha_j}
 =\sum_{k\ge0}(-1)^k h_k(\boldsymbol\alpha;\boldsymbol d)
                                      \zeta(A+k,c).
 \label{utj:eq:multizeta}
\end{equation}
The left series initially requires $\operatorname{Re}A>1$; the right
is its meromorphic continuation. The only possible hyperplanes are
$A=1-k$, and their residues are
$(-1)^k h_k|_{A=1-k}$. Normal convergence of the large-$k$ tail follows
from $|h_k|\le C R_0^k$ for $\max|d_j|<R_0<c$ and
$|\zeta(A+k,c)|\le Cc^{-k}$ on compact order sets.

For complex weights $w_j$, put $W=\sum_jw_j\ne0$ and
$\mathcal Z_{\boldsymbol w}(t)=\mathcal Z(t\boldsymbol w;\boldsymbol a)$.
Then the transverse ray has a removable singularity at zero and
\begin{align}
 \mathcal Z_{\boldsymbol w}(0)
   &=\frac12-\frac{\sum_jw_ja_j}{W},\label{utj:eq:rayvalue}\\
 \mathcal Z_{\boldsymbol w}'(0)
   &=\sum_jw_j\log\Gamma(a_j)-\frac W2\log(2\pi).
 \label{utj:eq:raygamma}
\end{align}
To prove this, isolate $k=0,1$ in \eqref{utj:eq:multizeta}. With
$D=\sum_jw_jd_j$, these terms are
$\zeta(Wt,c)-tD\zeta(1+Wt,c)$; the second has constant $-D/W$
and linear coefficient $-D\gamma_0(c)$. For $k\ge2$,
$h_k(0)=0$ and $\partial_t h_k(t\boldsymbol w;\boldsymbol d)|_0
 =\sum_jw_jd_j^k/k$. Summing the convergent Taylor series
\[
 \log\Gamma(c+d)-\log\Gamma(c)-d\psi(c)
    =\sum_{k\ge2}\frac{(-1)^kd^k}{k}\zeta(k,c)
\]
gives \eqref{utj:eq:raygamma}, using $\gamma_0(c)=-\psi(c)$.
The ray value generally depends on direction: for shifts $1,2$, the
rays with weights $(1,1)$ and $(2,1)$ give $-1$ and $-5/6$.
Thus a joint value is not being assigned at the multivariable origin.

The full continuation report also proves the unilateral multishift
Lerch expansion, all ordinary primitives, the meromorphic polar
hyperplanes, every transverse spectral-ray jet, and the pole-aware
cyclotomic Stieltjes filter. Its proofs and execution records retain
separate status: 824 finite exact assertions and 96 numerical comparisons
at each of two working precisions corroborate, but do not replace, the
analytic arguments. No global priority or proof-assistant claim is made.
